% Part: first-order-logic
% Chapter: syntax-and-semantics
% Section: satisfaction

\documentclass[../../../include/open-logic-section]{subfiles}

\begin{document}

\olfileid{fol}{syn}{sat}

\olsection{公式在结构中的满足}

\begin{explain}
联系项与公式等表达式和结构的基本概念，是项的\emph{值}与公式的\emph{满足}。
非形式地说，项的值是结构中的一个元素——如果项只是一个常项，则其值就是结构赋予该常项的对象；
如果项是用函数符号构造出来的，那么它的值就由各常项的值
以及结构指派给项中函数符号的函数计算得出。
一个公式在结构中被\emph{满足}，
是指结构对谓词的解释使得该公式在结构的论域中为真。
这种满足关系是归纳规定的：
结构的规定直接确定了何时原子公式被满足；而对于复合公式，我们则根据其主联结词或量词，
以及直接子公式是否被满足来定义其满足条件。

量词的情形则有点棘手，因为量词公式的直接子公式含有自由变元，
而结构本身并不指定变元的值。
为了处理这个困难，我们还要引入\emph{变元指派}，
并且，我们不是仅相对于结构来定义满足，而是相对于结构连同变元指派来定义。
\end{explain}

\begin{defn}[变元指派]
结构~$\Struct{M}$ 的\emph{变元指派}~$s$ 是一个函数，它将每个变元映射到~$\Domain M$ 的一个元素，即 $s\colon \Var \to \Domain M$。
\end{defn}

\begin{explain}
结构为每个常项符号赋予一个值，而变元指派则为每个变元赋予一个值。
但我们想要用由它们构造出来的项，来命名论域中的元素。
为此，我们归纳地定义项的值。
对于常项符号和变元，其值就是结构或变元指派所规定的值；
对于更复杂的项，则利用结构指派给函数符号的那些函数递归地计算出来。
\end{explain}

\begin{defn}[项的值]
若 $t$ 是语言~$\Lang L$ 的一个项，$\Struct M$ 是 $\Lang L$ 的一个结构，且 $s$ 是 $\Struct M$ 的一个变元指派，则 $t$ 的\emph{值}~$\Value{t}{M}[s]$ 定义如下：
\begin{enumerate}
\item \indcase{t}{c}{$\Value{\indfrm}{M}[s] = \Assign{\indcomplex}{M}$。}
\item \indcase{t}{x}{$\Value{\indfrm}{M}[s] = s(\indcomplex)$。}
\item \indcase{t}{\Atom{f}{t_1, \ldots, t_n}}{
\[
\Value{\indfrm}{M}[s] = \Assign{f}{M}(\Value{t_1}{M}[s], \ldots,
\Value{t_n}{M}[s]).
\]}
\end{enumerate}
\end{defn}

\begin{defn}[$x$-变体]
若 $s$ 是结构~$\Struct M$ 的一个变元指派，则 $\Struct M$ 的任何至多在对 $x$ 的赋值上与 $s$ 不同的变元指派 $s'$，都称为 $s$ 的一个\emph{$x$-变体}。
若 $s'$ 是 $s$ 的一个 $x$-变体，我们记作 $\varAssign{s'}{s}{x}$。
\end{defn}

\begin{explain}
注意，指派 $s$ 的一个 $x$-变体并\emph{不一定}给 $x$ 赋予不同的值。
事实上，每个指派都算作其自身的一个 $x$-变体。
\end{explain}

\begin{defn}
  若 $s$ 是结构~$\Struct M$ 的一个变元指派，且 $m \in \Domain{M}$，则指派 $\Subst{s}{m}{x}$ 定义为
  \[\Subst{s}{m}{x}(y) = \begin{cases}
    m & \text{若 } y \ident x\\
    s(y) & \text{否则}.
  \end{cases}\]
\end{defn}

换句话说，$\Subst{s}{m}{x}$ 是 $s$ 的那个特定的 $x$-变体，它将论域元素 $m$ 赋给 $x$，而对 $x$ 以外的变元则赋予与 $s$ 相同的值。

\begin{defn}[满足]
\ollabel{defn:satisfaction}
公式~$!A$ 在结构~$\Struct M$ 中相对于变元指派~$s$ 被满足，记作：
$\Sat{M}{!A}[s]$，递归定义如下。（我们用
$\Sat/{M}{!A}[s]$ 表示“并非 $\Sat{M}{!A}[s]$”。）
\begin{enumerate}
\tagitem{prvFalse}{%
  \indcase{!A}{\lfalse}{$\Sat/{M}{\indfrm}[s]$.}}{}

\tagitem{prvTrue}{%
  \indcase{!A}{\ltrue}{$\Sat{M}{\indfrm}[s]$.}}{}

\item \indcase{!A}{\Atom{R}{t_1, \dots, t_n}}{$\Sat{M}{\indfrm}[s]$
  当且仅当 $\langle \Value{t_1}{M}[s], \dots, \Value{t_n}{M}[s] \rangle \in
  \Assign{R}{M}$.}

\item \indcase{!A}{\eq[t_1][t_2]}{$\Sat{M}{\indfrm}[s]$ 当且仅当
  $\Value{t_1}{M}[s] = \Value{t_2}{M}[s]$.}

\tagitem{prvNot}{%
  \indcase{!A}{\lnot !B}{$\Sat{M}{\indfrm}[s]$ 当且仅当
    $\Sat/{M}{!B}[s]$.}}{}

\tagitem{prvAnd}{%
  \indcase{!A}{(!B \land !C)}{$\Sat{M}{\indfrm}[s]$ 当且仅当 $\Sat{M}{!B}[s]$
    且 $\Sat{M}{!C}[s]$。}}{}

\tagitem{prvOr}{%
  \indcase{!A}{(!B \lor !C)}{$\Sat{M}{\indfrm}[s]$ 当且仅当
    $\Sat{M}{!B}[s]$ 或 $\Sat{M}{!C}[s]$（或两者皆真）。}}{}

\tagitem{prvIf}{%
  \indcase{!A}{(!B \lif !C)}{$\Sat{M}{\indfrm}[s]$ 当且仅当 $\Sat/{M}{!B}[s]$
    或 $\Sat{M}{!C}[s]$（或两者皆真）。}}{}

\tagitem{prvIff}{%
  \indcase{!A}{(!B \liff !C)}{$\Sat{M}{\indfrm}[s]$ 当且仅当要么
    $\Sat{M}{!B}[s]$ 与 $\Sat{M}{!C}[s]$ 均成立，要么 $\Sat{M}{!B}[s]$
    与 $\Sat{M}{!C}[s]$ 均不成立。}}{}

\tagitem{prvAll}{%
  \indcase{!A}{\lforall[x][!B]}{$\Sat{M}{\indfrm}[s]$ 当且仅当对于每个
    元素~$m \in \Domain M$，都有 $\Sat{M}{!B}[\Subst{s}{m}{x}]$。}}{}

\tagitem{prvEx}{%
  \indcase{!A}{\lexists[x][!B]}{$\Sat{M}{\indfrm}[s]$ 当且仅当至少存在一个
  元素~$m \in \Domain M$，使得 $\Sat{M}{!B}[\Subst{s}{m}{x}]$。}}{}
\end{enumerate}
\end{defn}

\begin{explain}
变元指派在最后\iftag{notprvEx,notprvAll}{一个子句}{两个子句}中至关重要。\iftag{prvAll}{我们不能通过“对于所有 $m \in \Domain{M}$，$\Sat{M}{!B(m)}$”来定义 $\lforall[x][!B(x)]$ 的满足。}{}\iftag{prvEx}{我们不能通过“至少存在一个 $m \in \Domain{M}$，$\Sat{M}{!B(m)}$”来定义 $\lexists[x][!B(x)]$ 的满足。}{} 原因是，如果 $m \in \Domain M$，它并不是语言中的符号，因此 $!B(m)$~不是一个公式（也就是说，$\Subst{!B}{m}{x}$ 是未定义的）。我们也不能假定有常项符号或项能够为 $\Struct{M}$ 的每个元素命名，因为结构的定义中并没有这一要求。在标准语言中，常项符号的集合是可数无穷的，因此如果 $\Domain{M}$ 是不可数的，那么甚至连常项符号都不足以命名每个对象。

我们通过引入变元指派来解决这个问题，它允许我们将变元直接与论域中的元素联系起来。于是，我们不说例如 \iftag{prvEx}{$\lexists[x][!B(x)]$}{$\lforall[x][!B(x)]$} 在~$\Struct M$ 中被满足当且仅当 \iftag{prvEx}{至少存在一个 $m \in \Domain{M}$}{对于所有 $m \in \Domain{M}$，$\Sat{M}{!B(m)}$}，而是说它在~$\Struct M$ 中\emph{相对于}~$s$ 被满足，当且仅当对于\iftag{prvEx}{至少一个}{每一个}$m \in \Domain M$，$!B(x)$ 相对于~$\Subst{s}{m}{x}$ 被满足。
\end{explain}

\begin{ex}
令语言为 $\Lang{L} = \{a, b, f, R\}$，其中 $a$ 和 $b$ 是常项符号，$f$~是一个二元函数符号，$R$~是一个二元谓词符号。
考虑结构~$\Struct{M}$，其定义为：
\begin{enumerate}
\item $\Domain M = \{1, 2, 3, 4\}$
\item $\Assign{a}{M} = 1$
\item $\Assign{b}{M} = 2$
\item $\Assign{f}{M}(x, y) = x+y$ 如果 $x+y \le 3$，否则为 $3$。
\item $\Assign{R}{M} = \{\tuple{1, 1}, \tuple{1, 2}, \tuple{2, 3}, \tuple{2, 4}\}$
\end{enumerate}
函数 $s(x) = 1$ 将 $1 \in \Domain{M}$ 赋予每个变元，它是~$\Struct{M}$ 的一个变元指派。

那么，
\begin{align*}
\Value{f(a,b)}{M}[s] & = \Assign{f}{M}(\Value{a}{M}[s], \Value{b}{M}[s]).
\intertext{由于 $a$ 和 $b$ 是常项符号，$\Value{a}{M}[s]
  = \Assign{a}{M} = 1$，且 $\Value{b}{M}[s] = \Assign{b}{M} = 2$。所以}
\Value{f(a,b)}{M}[s] & = \Assign{f}{M}(1, 2) = 1+2 = 3.
\intertext{为计算 $f(f(a,b),a)$ 的值，我们必须考虑}
\Value{f(f(a,b),a)}{M}[s] & = \Assign{f}{M}(\Value{f(a, b)}{M}[s],
\Value{a}{M}[s]) = \Assign{f}{M}(3, 1) = 3,
\intertext{因为 $3+1 > 3$。由于 $s(x) = 1$ 且 $\Value{x}{M}[s] =
  s(x)$，我们也有}
\Value{f(f(a,b),x)}{M}[s] & = \Assign{f}{M}(\Value{f(a, b)}{M}[s],
\Value{x}{M}[s]) = \Assign{f}{M}(3, 1) = 3,
\end{align*}

一个原子公式~$R(t_1, t_2)$ 被满足，当且仅当其自变元的值构成的元组，即 $\tuple{\Value{t_1}{M}[s],
  \Value{t_2}{M}[s]}$，是~$\Assign{R}{M}$ 的一个元素。所以，例如，我们有 $\Sat{M}{R(b,f(a,b))}[s]$，因为 $\tuple{\Value{b}{M},
  \Value{f(a,b)}{M}} = \tuple{2, 3} \in \Assign{R}{M}$；但 $\Sat/{M}{R(x, f(a,b))}[s]$，因为 $\tuple{1, 3} \notin \Assign{R}{M}[s]$。

要判断一个非原子公式~$!A$ 是否被满足，需根据其主联结词应用归纳定义中的相应条款。
例如，$R(a, a) \lif (R(b,
x) \lor R(x, b))$ 的主联结词是~$\lif$，并且
\begin{align*}
  & \Sat{M}{R(a, a) \lif (R(b, x) \lor R(x, b))}[s] \text{ 当且仅当 }\\
  & \qquad
  \Sat/{M}{R(a,a)}[s] \text{ 或 } \Sat{M}{R(b, x) \lor R(x, b)}[s]
  \intertext{由于 $\Sat{M}{R(a,a)}[s]$（因为 $\tuple{1,1} \in
    \Assign{R}{M}$），我们还不能确定答案，必须首先弄清楚 $\Sat{M}{R(b, x) \lor R(x, b)}[s]$ 是否成立：}
  & \Sat{M}{R(b, x) \lor R(x, b)}[s] \text{ 当且仅当 }\\
  & \qquad \Sat{M}{R(b,x)}[s] \text{ 或 } \Sat{M}{R(x, b)}[s]
  \intertext{而这是成立的，因为 $\Sat{M}{R(x, b)}[s]$（基于 $\tuple{1,2} \in \Assign{R}{M}$）。}
\end{align*}

回顾 $s$ 的一个 $x$-变体，是指一个至多在对 $x$ 的赋值上与 $s$ 不同的变元指派。
对于~$\Domain{M}$ 中的每个元素，都存在 $s$ 的一个 $x$-变体：
\begin{align*}
  s_1 & = \Subst{s}{1}{x}, &
  s_2 & = \Subst{s}{2}{x},\\
  s_3 & = \Subst{s}{3}{x}, &
  s_4 & = \Subst{s}{4}{x}.
\end{align*}
因此，例如，$s_2(x) = 2$，并且对于除~$x$ 外的所有变元~$y$，$s_2(y) = s(y) = 1$。
由于 $\Domain{M} = \{1, 2, 3, 4\}$，这些就是结构~$\Struct{M}$ 下 $s$ 的所有 $x$-变体。
注意，$s_1 = s$（$s$~总是其自身的 $x$-变体）。

\iftag{prvEx}{要判断一个存在量化公式~$\lexists[x][!A(x)]$ 是否被满足，我们必须确定是否存在至少一个 $m \in \Domain M$，使得 $\Sat{M}{!A(x)}[\Subst{s}{m}{x}]$。所以，
  \[
  \Sat{M}{\lexists[x][(R(b,x) \lor R(x,b))]}[s],
  \]
  因为 $\Sat{M}{R(b,x) \lor R(x, b)}[\Subst{s}{1}{x}]$
  （$\Subst{s}{3}{x}$ 也符合要求）。但是，
  \[
  \Sat/{M}{\lexists[x][(R(b,x) \land R(x,b))]}[s]
  \]
  因为无论选择哪个 $m \in \Domain{M}$，都有 $\Sat/{M}{R(b,x) \land R(x,b)}[\Subst{s}{m}{x}]$。}{}

\iftag{prvAll}{要判断一个全称量化公式~$\lforall[x][!A(x)]$ 是否被满足，我们必须确定是否对所有 $m \in \Domain M$ 都有 $\Sat{M}{!A(x)}[\Subst{s}{m}{x}]$。所以，
  \[
  \Sat{M}{\lforall[x][(R(x,a) \lif R(a,x))]}[s],
  \]
  因为对于所有 $m \in \Domain M$，$\Sat{M}{R(x,a) \lif R(a,x)}[\Subst{s}{m}{x}]$。对于 $m = 1$，我们有 $\Sat{M}{R(a,x)}[\Subst{s}{1}{x}]$，所以后件为真；对于 $m = 2$、$3$ 和~$4$，我们有 $\Sat/{M}{R(x,a)}[\Subst{s}{m}{x}]$，所以前件为假。但是，
  \[
  \Sat/{M}{\lforall[x][(R(a,x) \lif R(x,a))]}[s]
  \]
  因为 $\Sat/{M}{R(a,x) \lif R(x,a)}[\Subst{s}{2}{x}]$（由于 $\Sat{M}{R(a, x)}[\Subst{s}{2}{x}]$ 和 $\Sat/{M}{R(x,
  a)}[\Subst{s}{2}{x}]$）。}{}

\iftag{defEx}{要判断一个存在量化公式~$\lexists[x][!A(x)]$ 是否被满足，我们必须确定是否 $\Sat{M}{\lnot \lforall[x][\lnot !A(x)]}[s]$。例如，我们有
  \[
  \Sat{M}{\lexists[x][(R(b,x) \lor R(x,b))]}[s].
  \]
  首先，$\Sat{M}{R(b,x) \lor R(x, b)}[\Subst{s}{1}{x}]$
  （$\Subst{s}{3}{x}$ 也符合要求）。所以，$\Sat/{M}{\lnot(R(b,x) \lor R(x,b))}[\Subst{s}{1}{x}]$，因此 $\Sat/{M}{\lforall[x][\lnot((R(b,x) \lor R(x,b))]}[s]$，进而 $\Sat{M}{\lnot\lforall[x][\lnot((R(b,x) \lor
  R(x,b))]}[s]$。另一方面，
  \[
  \Sat/{M}{\lexists[x][(R(b,x) \land R(x,b))],}[s].
  \]
  这是因为 $\Sat{M}{\lforall[x][\lnot(R(b,x) \land
  R(x,b))]}[s]$，即不存在 $m \in \Domain M$ 使得 $\Sat{M}{R(b,x) \land
  R(x,b)}[\Subst{s}{m}{x}]$。从这些例子中或许可以猜到，$\Sat{M}{\lexists[x][!A(x)]}[s]$ 当且仅当至少存在一个 $m \in \Domain M$，使得 $\Sat{M}{!A(x)}[\Subst{s}{m}{x}]$。}{}

\iftag{defAll}{要判断一个全称量化公式~$\lforall[x][!A(x)]$ 是否被满足，我们必须确定是否 $\Sat{M}{\lnot\lexists[x][\lnot !A(x)]}[s]$。例如，
  \[
  \Sat{M}{\lforall[x][(R(x,a) \lif R(a,x))]}[s],
  \]
  首先，对于所有 $m \in \Domain M$，$\Sat{M}{R(x,a) \lif R(a,x)}[\Subst{s}{m}{x}]$（$m = 1$ 时 $\Sat{M}{R(a,x)}[\Subst{s}{1}{x}]$，$m = 2$、$3$ 或~$4$ 时 $\Sat/{M}{R(a,x)}[\Subst{s}{m}{x}]$）。
  因此，不存在 $m \in \Domain M$ 使得 $\Sat{M}{\lnot(R(x,a) \lif R(a,x))}[\Subst{s}{m}{x}]$，进而 $\Sat/{M}{\lexists[x][\lnot(R(x,a) \lif R(a,x))]}[s]$。所以，$\Sat{M}{\lnot\lexists[x][\lnot(R(x,a) \lif R(a,x))]}[s]$。另一方面，
  \[
  \Sat/{M}{\lforall[x][(R(a,x) \lif R(x,a))]}[s],
  \]
  因为 $\Sat/{M}{R(a,x) \lif R(x,a)}[\Subst{s}{2}{x}]$，所以 $\Sat{M}{\lexists[x][\lnot(R(a,x) \lif R(x,a))]}[s]$。从这些例子中或许可以猜到，$\Sat{M}{\lforall[x][!A(x)]}[s]$ 当且仅当对于每个 $m \in \Domain M$，都有 $\Sat{M}{!A(x)}[\Subst{s}{m}{x}]$。}{}

对于一个更复杂的情况，考虑：
\[
\lforall[x][(R(a,x) \lif \lexists[y][R(x,y)])].
\]
由于 $\Sat/{M}{R(a,x)}[\Subst{s}{3}{x}]$ 且 $\Sat/{M}{R(a,x)}[\Subst{s}{4}{x}]$，需要考虑条件句后件是否成立的有趣情形仅有 $m = 1$ 和 $m = 2$。$\Sat{M}{\lexists[y][R(x,y)]}[\Subst{s}{1}{x}]$ 是否成立？如果存在至少一个 $n \in \Domain M$ 使得 $\Sat{M}{R(x,y)}[\Subst{\Subst{s}{1}{x}}{n}{y}]$，则其成立。事实上，如果我们取 $n = 1$，则有 $\Subst{\Subst{s}{1}{x}}{n}{y} = \Subst{s}{1}{y} =
s$。由于 $s(x) = 1$，$s(y) = 1$，且 $\tuple{1,1} \in \Assign{R}{M}$，答案是肯定的。 

要判断 $\Sat{M}{\lexists[y][R(x,y)]}[\Subst{s}{2}{x}]$ 是否成立，我们必须考察变元指派 $\Subst{\Subst{s}{2}{x}}{n}{y}$。这里，对于 $n = 1$，该指派为~$s_2 = \Subst{s}{2}{x}$，它不满足 $R(x,y)$（$s_2(x) = 2$，$s_2(y) = 1$，且 $\tuple{2,1}\notin \Assign{R}{M}$）。然而，考虑 $\Subst{\Subst{s}{2}{x}}{3}{y} = \Subst{s_2}{3}{y}$。
$\Sat{M}{R(x,y)}[\Subst{s_2}{3}{y}]$，因为 $\tuple{2,3} \in
\Assign{R}{M}$，所以 $\Sat{M}{\lexists[y][R(x,y)]}[s_2]$。

所以，对所有 $n \in \Domain M$，要么 $\Sat/{M}{R(a,x)}[\Subst{s}{m}{x}]$（若 $m = 3$，$4$），要么 $\Sat{M}{\lexists[y][R(x,y)]}[\Subst{s}{m}{x}]$（若 $m = 1$，$2$），于是
\[
\Sat{M}{\lforall[x][(R(a,x) \lif \lexists[y][R(x,y)])]}[s].
\]
另一方面，
\[
\Sat/{M}{\lexists[x][(R(a,x) \land \lforall[y][R(x,y)])]}[s].
\]
只有当 $m = 1$ 和 $m = 2$ 时，我们才有 $\Sat{M}{R(a,x)}[\Subst{s}{m}{x}]$。但对于这两个 $m$ 值，又分别存在一个 $n \in \Domain M$，即 $n = 4$，使得 $\Sat/{M}{R(x,y)}[\Subst{\Subst{s}{m}{x}}{n}{y}]$，从而对于 $m = 1$ 和 $m = 2$，$\Sat/{M}{\lforall[y][R(x,y)]}[\Subst{s}{m}{x}]$。因此，不存在 $m \in \Domain M$ 使得 $\Sat{M}{R(a,x)
\land \lforall[y][R(x,y)]}[\Subst{s}{m}{x}]$。
\end{ex}

\iftag{defEx}{%


\begin{prop}\ollabel{prop:sat-ex}
$\Sat{M}{\lexists[x][!B(x)]}[s]$ 当且仅当存在 $s$ 的一个 $x$-变体 $s'$，使得 $\Sat{M}{!B(x)}[s']$。
\end{prop}

\begin{proof}
  练习。
\end{proof}


}{}

\tagprob{defEx}


\begin{prob}
  证明 \olref[fol][syn][sat]{prop:sat-ex}
\end{prob}


\tagendprob

\iftag{defAll}{%


\begin{prop}\ollabel{prop:sat-all}
$\Sat{M}{\lforall[x][!B(x)]}[s]$ 当且仅当对于 $s$ 的每个 $x$-变体~$s'$，都有
  $\Sat{M}{!B(x)}[s']$
\end{prop}

\begin{proof}
  练习。
\end{proof}


}{}

\tagprob{defAll}


\begin{prob}
  证明 \olref[fol][syn][sat]{prop:sat-all}
\end{prob}


\tagendprob

\begin{prob}
设 $\Lang L = \{c, f, A\}$ 是带有一个常项符号、一个一元函数符号和一个二元谓词符号的语言，并设结构~$\Struct{M}$ 定义如下：
\begin{enumerate}
\item $\Domain M = \{1, 2, 3\}$
\item $\Assign{c}{M} = 3$
\item $\Assign{f}{M}(1) = 2, \Assign{f}{M}(2) = 3, \Assign{f}{M}(3) = 2$
\item $\Assign{A}{M} = \{\tuple{1, 2}, \tuple{2, 3}, \tuple{3, 3}\}$
\end{enumerate}
(a) 设对所有变元~$v$ 有 $s(v) = 1$。判断 \[
\Sat{M}{\lexists[x][(A(f(z), c) \lif \lforall[y][(A(y, x) \lor A(f(y),
      x))])]}[s]
\] 是否成立。
解释原因。

(b) 给出另一个结构和变元指派，使得该公式不被满足。
\end{prob}

\end{document}